\section{The studio GUI}

The GUI (\cmd{gui/app.py}, Streamlit) presents the pipeline as a
gated workflow across tabs: Specimen \& Fabric, Probe \& Excitation,
Error model, Acquisition, then Backscatter and COF reconstruction,
the last two unlocking only once acquired data exists in the
session. The FW sweep tab (\cmd{gui/sweep\_tab.py}) configures,
launches and monitors full-wave sweeps through the sweep engine.

The specimen views are worth knowing about even for command-line
users, since they are the fastest way to sanity-check a fabric: a 3D
per-grain isosurface rendering coloured by c-axis, a c-axis vector
field, an equal-area pole figure with orientation-tensor
eigenvalues, and 2D slices of the grain and azimuth fields.

One engineering note recorded here because it is invisible in the
code's behaviour and expensive to rediscover: the GUI never probes
the GPU at import time. Streamlit executes the script in a worker
thread, and initialising CUDA off the main thread crashes the whole
server without a traceback; the backend is therefore detected by
inspection (is cupy importable, is a CUDA path configured) and the
solvers touch the driver only when actually run.
